\appendixchapter{Hints for Selected Exercises}
\label{app:hints}
These optional hints suggest a starting point, not a completed solution.
Try the exercise first; return to its hypotheses after reading the hint.
The exercise number in each entry links back to the original problem.
Numbers restart within each chapter or appendix, exactly as in the exercise
sets. A superscript asterisk means \emph{hint available}, not ``hard exercise.''
The optional appendices have more selected hints to support independent reading.

\par\Needspace{9\baselineskip}
\section*{Chapter 2}
\addcontentsline{toc}{section}{Chapter 2}

\exercisehint{strings}{Group strings by their length. Each group is finite; arrange the groups by increasing length and list each group in a fixed order.}

\exercisehint{images}{For each inclusion, start with an arbitrary element and unpack image and inverse image separately. For failure of equality, try a constant function on a two-element set.}

\exercisehint{induction-sum}{Write the sum through $n+1$ as the sum through $n$ plus its last term. Use the induction hypothesis only on the first part.}

\par\Needspace{9\baselineskip}
\section*{Chapter 3}
\addcontentsline{toc}{section}{Chapter 3}

\exercisehint{representatives}{Compare a constant rational sequence with that constant plus $1/n$. Equivalence concerns the limit of a difference, not equality of individual terms.}

\exercisehint{rational-approximations}{Insert $x_j$ and $x_k$ between $r_j$ and $r_k$ in the triangle inequality. Give each of the three errors a share of the tolerance.}

\par\Needspace{9\baselineskip}
\section*{Chapter 4}
\addcontentsline{toc}{section}{Chapter 4}

\exercisehint{seq-direct}{Simplify the error first. Choose a natural number strictly larger than your sufficient bound; verify the estimate for every $n\geq N$.}

\exercisehint{seq-laws}{Divide by the leading powers and bound the normalized denominator away from zero before applying a quotient law.}

\exercisehint{seq-geometric}{Use Bernoulli on $4/3$ for decay and on $3/2$ for growth. For the negative base separate odd and even indices.}

\exercisehint{seq-radical}{Multiply by the conjugate, then divide the denominator by $k$.}

\exercisehint{seq-model}{Let $e_n$ be the error after $n$ corrections, starting at $n=0$. Inductively bound $|e_n|$ by a geometric sequence; equality is not assumed.}

\exercisehint{seq-disprove}{For the first use one fixed positive error gap. For the second, either exclude every candidate or exhibit two distinct subsequential limits.}

\exercisehint{supremum-proof}{The number $s-\eps$ cannot be an upper bound. Find one term above it and use monotonicity for all subsequent indices.}

\exercisehint{decreasing-completeness}{Let the candidate be the infimum of the range. What does the definition of infimum say about a term below that candidate plus $\eps$?}

\par\Needspace{9\baselineskip}
\section*{Chapter 5}
\addcontentsline{toc}{section}{Chapter 5}

\exercisehint{fun-direct}{For the square first restrict $|x-3|<1$; for the reciprocal keep $|x+1|$ above a fixed positive number.}

\exercisehint{fun-infinity}{Normalize by the degrees; for the negative end set $x=-t$. Retain the parity of the degree difference.}

\exercisehint{fun-removable}{Compute the punctured limits before examining the assigned values.}

\exercisehint{fun-pairs}{For positive results estimate differences. For negative results use input pairs approaching each other while the outputs retain a fixed gap.}

\exercisehint{fun-ivt-model}{Continuity and the endpoint readings give existence. No monotonicity assumption has been supplied.}

\exercisehint{fun-sensitivity}{Factor the output difference and bound the sum of the two inputs over the entire domain.}

\exercisehint{jump-continuity}{At the junction compare the sequences $1/n$ and $-1/n$ and the assigned value. Away from the junction choose a neighborhood lying in one branch.}

\par\Needspace{9\baselineskip}
\section*{Chapter 6}
\addcontentsline{toc}{section}{Chapter 6}

\exercisehint{uniform-cases}{Inspect $-1,0,1$ separately. On $[0,r]$ use $r^n$; near $1$ use moving inputs.}

\exercisehint{uniform-model}{Complete the square in $x(1-x)$ to locate its largest possible value.}

\exercisehint{moving-spike}{For nonuniformity choose the input so that $nx$ stays equal to $1$. On $[\delta,1]$ use $1+n^2x^2\geq n^2x^2$.}

\par\Needspace{9\baselineskip}
\section*{Chapter 7}
\addcontentsline{toc}{section}{Chapter 7}

\exercisehint{series-selection}{Start with the term test; then look for cancellation, a power benchmark, a ratio, or an $n$th root, in that order.}

\exercisehint{series-model}{Index the first contribution carefully. Apply the finite geometric identity before passing to the limiting total.}

\exercisehint{series-preview}{Write finite upper bounds $b$ on the integrals, evaluate, and then let $b\to\infty$. Keep this calculation separate from the condensation proof.}

\exercisehint{three-endpoints}{Test each endpoint independently. For uniform convergence on the closed interval, a summable bound independent of $x$ is useful; at a divergent endpoint uniform convergence is impossible.}

\par\Needspace{9\baselineskip}
\section*{Chapter 8}
\addcontentsline{toc}{section}{Chapter 8}

\exercisehint{derivative-definition}{Keep $|h|<|a|/2$ in the reciprocal quotient at $a\ne0$. For the piecewise function compare the two one-sided quotients.}

\exercisehint{optimization-model}{Use the perimeter condition to express one side in terms of the other. Compare the interior critical point with the boundary values.}

\exercisehint{log-taylor}{Bound the third derivative on the whole interval from $0$ to $0.1$, then use the Lagrange remainder. A decimal guess is not an error bound.}

\par\Needspace{9\baselineskip}
\section*{Chapter 9}
\addcontentsline{toc}{section}{Chapter 9}

\exercisehint{darboux-direct}{For the linear function, sum the oscillation on each cell times its width. For the other function use density in every cell.}

\exercisehint{flow-model}{Accumulate the given rate over elapsed time. Divide the total accumulation by the duration to obtain its average rate.}

\exercisehint{integral-spikes}{Place a triangle on $[1/(2n),1/n]$ and choose its height to keep its area fixed. Check the value at zero separately.}

\par\Needspace{9\baselineskip}
\section*{Appendix A}
\addcontentsline{toc}{section}{Appendix A}

\exercisehint{app-a-1}{Fix $a$ and recurse from $1$ by multiplication by $a$. For the addition law, fix $m$ and induct on $n$.}

\exercisehint{app-a-3}{If $n$ is not prime, choose a proper divisor between $2$ and $n-1$ and apply the induction hypothesis to that divisor.}

\exercisehint{app-a-7}{Order the independent sets containing the prescribed set by inclusion. For a chain, test a dependence in its union using only the finitely many vectors involved.}


\par\Needspace{9\baselineskip}
\section*{Appendix B}
\addcontentsline{toc}{section}{Appendix B}

\exercisehint{app-b-2}{In a poset category, arrows in both directions force equality. Apply this to divisibility, then use transitivity rather than inventing underlying functions.}

\exercisehint{app-b-3}{Reverse the displayed arrows first: in $\mathcal C$ they have the form $C\to B\to A$. The composite in the opposite category is the formal opposite of their composite in $\mathcal C$.}

\exercisehint{app-b-5}{Evaluate each proposed composite on a covector and then on a vector. Precomposition determines the reversed arrow order.}

\exercisehint{app-b-6}{Write both composites explicitly. One route applies $f$ before taking the diagonal; the other applies $f$ to each coordinate after taking it.}

\exercisehint{app-b-8}{Start with $\omega\in\Omega^k(N)$. Both routes must end in $\Omega^{k+1}(M)$; this requirement fixes every vertical arrow direction.}


\par\Needspace{9\baselineskip}
\section*{Appendix C}
\addcontentsline{toc}{section}{Appendix C}

\exercisehint{app-c-2}{The $(i,j)$ entry of the matrix product is a cofactor expansion. For $i\ne j$, identify a matrix with two equal rows.}

\exercisehint{app-c-5}{Every pair of positions is an inversion in the reversing list. Count pairs before evaluating the sign.}

\exercisehint{app-c-6}{Reduce the integer determinant in the field. The field equality $-1=1$ does not change the parity of a permutation.}


\par\Needspace{9\baselineskip}
\section*{Appendix D}
\addcontentsline{toc}{section}{Appendix D}

\exercisehint{app-d-1}{A linear map descends exactly when it vanishes on the relation subspace. Evaluate both candidates on $(1,1)$, including in characteristic two.}

\exercisehint{app-d-2}{Check that $\pi_H\circ L$ kills $R$, then invoke the quotient universal property.}

\exercisehint{app-d-5}{Compare the values of the proposed bilinear map on $(e_1,e_2)$ and $(e_2,e_1)$.}

\exercisehint{app-d-7}{Project to the indicated two-dimensional factors. An elementary tensor has a coefficient table whose columns are proportional.}

\exercisehint{app-d-11}{First test $L_{\mathbf v}$ on elementary wedges in the remaining slots; then use that those wedges span.}


\par\Needspace{9\baselineskip}
\section*{Appendix E}
\addcontentsline{toc}{section}{Appendix E}

\exercisehint{app-e-1}{First prove constancy on integer translates. For the topology, factor through the compact interval model and use the compact-to-Hausdorff criterion.}

\exercisehint{app-e-3}{Use a signed horizontal coordinate on the two sides of the seam. The negative side must apply the edge reversal.}

\exercisehint{app-e-6}{Identify the map with angle doubling. After checking its fibers, use compactness and the Hausdorff target.}


\par\Needspace{9\baselineskip}
\section*{Appendix F}
\addcontentsline{toc}{section}{Appendix F}

\exercisehint{app-f-1}{Try $(x,y,z)\mapsto(x,y,z-f(x,y))$; its inverse adds $f$ back.}

\exercisehint{app-f-2}{Compute $Dh$ and inspect the two axes in the zero fiber near their crossing.}

\exercisehint{app-f-3}{Restrict $dz$ to the tangent plane of the sphere. At a regular point, intersect the two tangent kernels.}

\exercisehint{app-f-4}{Use the sequence approaching the omitted parameter endpoint. For the compact case, combine immersion with the compact-to-Hausdorff criterion.}

\exercisehint{app-f-5}{Differentiate $A^TA$ with symmetric matrices as target. To test surjectivity at an orthonormal frame, try a variation of the form $H=AB$ with $B$ symmetric.}

\exercisehint{app-f-7}{Take products of adapted charts and reorder the coordinates so that the free coordinates come first.}


\par\Needspace{9\baselineskip}
\section*{Appendix G}
\addcontentsline{toc}{section}{Appendix G}

\exercisehint{app-g-1}{Differentiate $x^2+y^2=1$ at $(1,0)$ and then apply the derivative of left multiplication.}

\exercisehint{app-g-2}{Left multiplication by $A$ is a diffeomorphism. Translate the already computed identity tangent space.}

\exercisehint{app-g-3}{Compute $[X,Y]$ on the coordinate functions first; these determine its coefficients.}

\exercisehint{app-g-4}{Multiply the displayed matrices in both orders. Use bilinearity and antisymmetry to reduce repeated work.}

\exercisehint{app-g-5}{For the diagonal matrix compute powers entrywise; for the other matrix use $A^2=0$ to truncate the series.}

\exercisehint{app-g-6}{In the $SO(2)$ realization, the initial tangent is $aJ$. Apply the one-parameter subgroup result and $J^2=-I$.}


\clearpage
\section*{Appendix H}
\addcontentsline{toc}{section}{Appendix H}

\exercisehint{app-h-1}{For $P\,dx+Q\,dy$, compare $\partial_xQ$ with $\partial_yP$. A primitive must have these two partial derivatives.}

\exercisehint{app-h-2}{For $k=2$, evaluate the two-form on $(x,v)$ and include the factor $t$ before integrating.}

\exercisehint{app-h-3}{Identify the interval components. A smooth locally constant function may take a nonzero value on every component, so use a product rather than a direct sum.}

\exercisehint{app-h-4}{Pull back the angular form along $t\mapsto mt$ and compare its period. In degree zero use connectedness.}

\exercisehint{app-h-5}{Compute closedness directly, then integrate around a circle in the annulus. Use the local Poincar\'e lemma for the last claim.}

\exercisehint{app-h-6}{Pull back the two degree-one generators separately. Record their periods on the two coordinate circles before writing the matrix.}

\exercisehint{app-h-7}{Use the product parametrization and orientation. A nonzero integral rules out exactness, but does not show that every other closed two-form is a multiple modulo exact forms.}

